%% kapa: exam=2 round=32 year=2021 period=3 subject=compensation_law subject_ko=감정평가 및 보상법규 seq=2 number=2 points=30
%% source: Q-Net 감정평가사 기출문제 (cst00309) artlSeq=5208993 file=raw/1008/5208993/감평 2차 시험지(2021).pdf
%% conversion: kordoc 4.24.51 -> tools/gen2.py (auto); review=proofread-sonnet; flags=none
\begin{longquestion}{【문제 2】}
甲은 A시에 토지를 소유하고 있다. A시장은 甲의 토지 등의 비교표준지로 A시 소재 일정 토지(2020. 1. 1. 기준 공시지가는 1 m\textsuperscript{2}당 1,000만원이다)를 선정하고, 甲의 토지 등과 비교표준지의 토지가격비준표상 총 가격배율을 1.00으로 조사함에 따라 甲의 토지의 가격을 1 m\textsuperscript{2}당 1,000만원으로 산정하였다. A시장으로부터 산정된 가격의 검증을 의뢰받은 감정평가사 乙은 甲의 토지가 비교표준지와 비교하여 환경조건, 획지조건 및 기타조건에서 열세에 있고, 특히 기타조건과 관련하여 비교표준지는 개발을 위한 거래가 이어지고 있으나 甲의 토지 등은 개발 움직임이 없다는 점을 ‘장래의 동향’으로 반영하여 91\%의 비율로 열세에 있다고 보아, 비교표준지의 공시지가를 약 83.9\%의 비율로 감액한 1 m\textsuperscript{2}당 839만원을 개별공시지가로 정함이 적정하다는 검증의견을 제시하였다. A시장은 A시 부동산가격공시위원회의 심의를 거쳐 이 검증의견을 그대로 받아들여 2020. 5. 20. 甲의 토지의 개별공시지가를 1 m\textsuperscript{2}당 839만원으로 결정·공시하고, 甲에게 개별 통지하였다. 甲은 토지가격비준표에 제시된 토지특성에 기초한 가격배율을 무시하고 乙이 감정평가방식에 따라 독자적으로 지가를 산정하여 제시한 검증의견을 그대로 반영하여 개별공시지가를 결정한 것은 위법하다고 보아, 「부동산 가격공시에 관한 법률」 제11조에 따라 2020. 6. 15. 이의신청을 제기하였고, 2020. 7. 10. 이의를 기각하는 내용의 이의신청결과가 甲에게 통지되었다. 다음 물음에 답하시오. (아래의 물음은 각 별개의 상황임) (30점) \par
\medskip\noindent 물음 1) 甲은 2020. 9. 10. 개별공시지가결정에 대해 취소소송을 제기하였다. 甲이 제기한 취소소송은 제소기간을 준수하였는가? (10점) \par
\medskip\noindent 물음 2) 甲이 개별공시지가결정에 대해 다투지 않은 채 제소기간이 도과하였고, 이후 甲의 토지에 대해 수용재결이 있었다. 甲이 보상금의 증액을 구하는 소송에서 개별공시지가결정의 위법을 주장하는 경우, 甲의 주장은 인용될 수 있는가? (20점) \par
\medskip
\begin{framedtext}
\textbf{참조 조문} \par
\textbf{<부동산 가격공시에 관한 법률>} \par
\textbf{제11조(개별공시지가에 대한 이의신청)} ① 개별공시지가에 이의가 있는 자는 그 결정·공시일부터 30일 이내에 서면으로 시장·군수 또는 구청장에게 이의를 신청할 수 있다. ② 시장·군수 또는 구청장은 제1항에 따라 이의신청 기간이 만료된 날부터 30일 이내에 이의신청을 심사하여 그 결과를 신청인에게 서면으로 통지하여야 한다. 이 경우 시장·군수 또는 구청장은 이의신청의 내용이 타당하다고 인정될 때에는 제10조에 따라 해당 개별공시지가를 조정하여 다시 결정·공시하여야 한다. \par
\textbf{<부동산 가격공시에 관한 법률 시행령>} \par
\textbf{제18조(개별공시지가의 검증)} ① <생략> \par
② 법 제10조제5항 본문에 따라 검증을 의뢰받은 감정평가법인등은 다음 각 호의 사항을 검토·확인하고 의견을 제시해야 한다. \par
\hspace*{1em}1. 비교표준지 선정의 적정성에 관한 사항 \par
\hspace*{1em}2. 개별토지 가격 산정의 적정성에 관한 사항 \par
\hspace*{1em}3. 산정한 개별토지가격과 표준지공시지가의 균형 유지에 관한 사항 \par
\hspace*{1em}4. 산정한 개별토지가격과 인근토지의 지가와의 균형 유지에 관한 사항 \par
\hspace*{1em}5. 표준주택가격, 개별주택가격, 비주거용 표준부동산가격 및 비주거용 개별부동산가격 산정 시 고려된 토지 특성과 일치하는지 여부 \par
\hspace*{1em}6. 개별토지가격 산정 시 적용된 용도지역, 토지이용상황 등 주요 특성이 공부(公簿)와 일치하는지 여부 \par
\hspace*{1em}7. 그 밖에 시장·군수 또는 구청장이 검토를 의뢰한 사항 \par
\textbf{<행정심판법>} \par
\textbf{제3조(행정심판의 대상)} ① 행정청의 처분 또는 부작위에 대하여 다른 법률에 특별한 규정이 있는 경우를 제외하고는 이 법에 의하여 행정심판을 제기할 수 있다. \par
\end{framedtext}
\end{longquestion}
